% 半人马小行星
% license CCBYSA3
% type Wiki

本文根据 CC-BY-SA 协议转载翻译自维基百科\href{https://en.wikipedia.org/wiki/Centaur_(small_Solar_System_body)}{相关文章}。

\begin{figure}[ht]
\centering
\includegraphics[width=6cm]{./figures/4207f1463529dc38.png}
\caption{} \label{fig_Centau_2}
\end{figure}
在行星天文学中，半人马小天体（centaur）是指一类在木星与海王星之间绕太阳运行**、并且**轨道会与一颗或多颗巨行星轨道相交的小型太阳系天体。由于这一特性，半人马天体的轨道通常是不稳定的；几乎所有已知半人马天体的轨道动力学寿命都只有数百万年**$^{[1]}$。不过，已知有一个例外——514107 Kaʻepaokaʻāwela，其轨道可能是稳定的（尽管是逆行轨道）$^{[2]}$（见注释$^{[1]}$）。

半人马天体通常同时表现出**小行星和彗星的特征**。它们的名称来源于神话中的“半人马”，即兼具马与人形态的生物。由于观测上对大尺寸天体存在偏向性，半人马天体的总体数量难以准确确定。对于直径大于 1 km 的半人马天体，其数量估计范围从**约 44,000 个**$^{[1]}$到**超过 1,000 万个**$^{[4][5]}$不等。

按照喷气推进实验室（JPL）的定义（也是本文采用的定义），**第一个被发现的半人马天体**是 **944 Hidalgo**，发现于 1920 年。然而，直到 1977 年发现 **2060 Chiron** 之后，半人马天体才被认识为一个独立的天体族群。迄今为止，已确认的最大半人马天体是 **10199 Chariklo**，其直径约为 **250 km**，大小相当于一颗中等尺寸的主带小行星，并且已知其拥有一套**环系统**。该天体发现于 1997 年。

目前尚无任何半人马天体被近距离拍摄过，不过有证据表明，**土星的卫星菲比（Phoebe）**——由卡西尼号探测器于 2004 年拍摄——可能是一颗来自**柯伊伯带**并被土星俘获的半人马天体$^{[6]}$。此外，**哈勃空间望远镜**也获取了一些关于 **8405 Asbolus** 表面特征的信息。

**谷神星（Ceres）**可能起源于外行星区域$^{[7]}$，如果这一假设成立，它可以被视为一颗“前半人马天体”（ex-centaur）；但目前观测到的半人马天体均起源于其他区域。

在已知占据半人马式轨道的天体中，约有 **30 个**显示出类似彗星的尘埃彗发，其中 **2060 Chiron**、**60558 Echeclus** 和 **29P/Schwassmann–Wachmann 1** 在**完全位于木星轨道之外**的情况下仍表现出可探测的挥发物释放$^{[8]}$。因此，Chiron 和 Echeclus 同时被归类为半人马天体和彗星，而 Schwassmann–Wachmann 1 一直被归类为彗星。其他半人马天体（如 **52872 Okyrhoe**）也被怀疑曾出现过彗发。理论上，任何被扰动并足够接近太阳的半人马天体，最终都可能演化为彗星。

---

### 分类

一颗天体若其**近日点距离或半长轴**位于外行星之间（即木星与海王星之间），即可被归类为半人马天体。由于该区域内轨道在长期尺度上的固有不稳定性，即便是当前并未与任何行星轨道相交的半人马天体（如 **2000 GM137** 和 **2001 XZ255**），其轨道也会逐渐演化，并最终受到扰动，从而开始与一颗或多颗巨行星的轨道相交$^{[1]}$。

一些天文学家仅将**半长轴位于外行星区域**的天体归为半人马天体；而另一些学者则认为，**只要近日点位于该区域内**即可归类为半人马天体，因为这两类轨道在动力学上同样是不稳定的。

---

如果你愿意，我可以继续帮你把**注释（note）部分**、**参考文献列表**或**与小行星带/柯伊伯带的动力学联系**统一整理成**完整 LaTeX 学术条目**。
